% ECONOMICS_PROBLEM: {"id": "J15-260", "title": "Reproducing racial wealth gaps", "jel_code": "J15", "source_id": "OP-0365"}
% !TeX program = xelatex
\documentclass[11pt,a4paper]{article}
\usepackage[margin=23mm]{geometry}
\usepackage{fontspec}
\setmainfont{Latin Modern Roman}
\usepackage{amsmath,amssymb,mathtools}
\usepackage{xcolor,enumitem,booktabs,longtable,array,xurl,hyperref}
\definecolor{muted}{HTML}{53636C}
\hypersetup{colorlinks=true,urlcolor=blue,linkcolor=blue}
\setlength{\parindent}{0pt}
\setlength{\parskip}{5pt}
\newcommand{\R}{\mathbb R}
\newcommand{\E}{\mathbb E}
\newcommand{\N}{\mathbb N}
\newcommand{\ind}{\mathbf 1}
\newcommand{\Var}{\operatorname{Var}}
\newcommand{\Cov}{\operatorname{Cov}}
\newcommand{\supp}{\operatorname{supp}}
\DeclareMathOperator*{\argmin}{arg\,min}
\DeclareMathOperator*{\argmax}{arg\,max}
\newcommand{\entrymeta}[3]{{\sffamily\small\color{muted}\textbf{\texttt{#1}}\quad|\quad #2\par #3\par}}
\newcommand{\Problemref}[1]{\texttt{#1}}
\newcommand{\JELref}[2]{\texttt{#2} (JEL \texttt{#1})}
\begin{document}
\section*{Reproducing racial wealth gaps}
\textbf{Problem ID:} \texttt{J15-260}\quad \textbf{JEL:} \texttt{J15}\par
% BEGIN ECONOMICS STATEMENT
\label{problem:OP-0365}
{\sffamily\small\color{muted}Original subject group: Inequality and economic mobility\par}
\label{definitions:problem:30007647}
\textbf{Audit classification:} Published research agenda. Authors, title, 2025 volume and DOI checked. The abstract explicitly calls for these research directions; policy causal contrasts are editorial.
\subsection*{Definitions and precise research target}
\textbf{Complete editorial problem.} Fix a United States cohort and define racial groups using a consistent self-identification rule. Let \(W_{it}\) be household net wealth, including business assets and debts, and let \(G_t=E[W_{it}\mid g_i=B]-E[W_{it}\mid g_i=W]\) be the Black--White mean wealth difference. Median and distributional gaps should be reported separately. Let policy \(p\) specify a feasible wealth transfer, access intervention, or institutional reform, and define \(W_{it}(p)\) allowing effects on subsequent education, employment, entrepreneurship, and transfers to descendants. Estimate \(\Delta G_t(p)=G_t(p)-G_t(0)\) over a declared multiyear horizon. The mechanism task is to distinguish wealth's consequences for later outcomes from the reverse effect of those outcomes on wealth, including community-level spillovers. Race itself is not the treatment. State assumptions needed to identify policy contrasts and any mediated pathway, accounting for selective household formation, survival, and business valuation. Determine which interventions durably change the distribution rather than simply shifting a single cross-sectional statistic. Historical decompositions and contemporary correlations provide evidence, but causal channel shares and policy transportability require further identification and explicit equilibrium assumptions.

\textbf{Definitions and admissible scope.} Fix whether units are people or households, and assign group membership and wealth consistently when households change. In a person-based baseline, net wealth \(W_{it}\) uses declared ownership shares, currency, date, assets, and debts; means require \(E|W_{it}|<\infty\). Define \(G_t(p)=E[W_t(p)\mid g=B]-E[W_t(p)\mid g=W]\) on the same baseline population and let \(\Delta G_t=G_t(p)-G_t(0)\). A reduction in this signed gap means movement toward zero only when its sign is known; use \(|G_t(p)|-|G_t(0)|\) for gap magnitude. Policies target resources or institutions, while race is a recorded group. Community effects require group or neighborhood assignment outcomes; a national wealth gap is not the sum of isolated household treatment effects. Use a fixed baseline household cohort and probability weights, with \(g=B\) denoting Black and \(g=W\) White households by the stated rule. Following household splits, assign assets and debts to baseline adult members through a stated ownership rule and reconstruct each original cohort unit, so postpolicy household formation does not silently change the target population. Define \(G_t(p)=E[W_{it}(p)\mid g_i=B]-E[W_{it}(p)\mid g_i=W]\) using that law and finite absolute means. A negative baseline gap improves in the mean-gap criterion when \(\Delta G_t>0\), while distributional improvement uses a separately declared criterion. Fiscal financing, survival, descendant links, and community exposure are parts of the complete intervention law.
\subsection*{Variants and assumption changes}
\begin{enumerate}
\item Wealth-to-opportunity direction (source-motivated): randomize or credibly identify wealth access and measure later education, employment, and entrepreneurship. This tests wealth as a cause rather than solely as an outcome.
\item Community-wealth direction (source-motivated): add collective assets and local ownership institutions, distinguishing their value from household assets to avoid double counting. Estimate spillovers and persistence after the initial intervention ends.
\end{enumerate}
\subsection*{References and source audit}
\textbf{Support and access.} Authors, title, 2025 volume and DOI checked. The abstract explicitly calls for these research directions; policy causal contrasts are editorial. \textbf{Locator:} Publisher abstract, paragraphs on wealth-to-outcome direction and community wealth.
\begin{enumerate}\raggedright
\item\hypertarget{definitions:source:30007647:1}{} Fenaba R. Addo; Daniel Auguste. 2025. \emph{Black and White Wealth Differentials in the United States: Explaining and Recreating Persistent Inequality}. Abstract, paragraphs 2–3.
\url{https://www.annualreviews.org/content/journals/10.1146/annurev-soc-090324-022010}.
DOI: \url{https://doi.org/10.1146/annurev-soc-090324-022010}.
\end{enumerate}

\subsection*{Common conventions: Source mathematical conventions}

These conventions apply to each entry unless explicitly replaced.

\textbf{Models and identification.} A primitive model \(m\) specifies all probability laws, feasible choices, information, equilibrium and measurement rules used by its target. The admissible class \(\mathcal M\) is fixed independently of outcomes. If \(P_O(m)\) is its observable probability law and \(\tau(m)\) the target, its identified set at observable law \(P\) is
\[
 \mathcal I_\tau(P)=\{\tau(m):m\in\mathcal M,\ P_O(m)=P\}.
\]
Point identification means that this set is a singleton. An empty set signals incompatibility of data and assumptions. Coordinate bounds need not describe the joint set. Bounds are sharp only if every retained target is attainable or arbitrarily approachable by compatible models. Finite bounds require explicit restrictions on unobserved quantities; unrestricted confounding, hidden wealth, extrapolation or arbitrary equilibrium selection can yield uninformative sets.

\textbf{Causal interventions.} A potential outcome \(Y_i(p)\) is the outcome under a complete policy regime \(p\), including timing, financing, information and market spillovers. The target population and units are fixed before treatment. With interference, use the complete assignment vector or a justified exposure mapping; individual no-interference notation alone cannot identify an economy-wide intervention. A difference of mean potential outcomes does not require the cross-world joint distribution. Conditioning on post-treatment migration, survival, enrollment or employment changes the estimand and needs separate justification.

\textbf{Statistical statements.} An estimator, confidence set, forecast or test specifies a sampling law, model class and loss/coverage criterion. Uniform \((1-\alpha)\)-coverage means \(\inf_{m\in\mathcal M}\Pr_m\{\tau(m)\in C_n\}\geq1-\alpha\), or an explicitly asymptotic version. The whole parameter space trivially has coverage, so informativeness requires a stated width, risk or power objective. A forecasting improvement is relative to a named benchmark, information set and evaluation population.

\textbf{Equilibrium and optimization.} An equilibrium comprises feasible plans, optimality, beliefs where needed, budget constraints and all market/resource clearing. The primitives or a stated benchmark define each correspondence \(\mathcal E(p;m)\). Nonemptiness is assumed or proved, never inferred just from compact policy sets. A selected-equilibrium target uses a declared selection rule; a robust target quantifies over all permitted equilibria. Use suprema or infima unless attainment is established. An \(\varepsilon\)-optimal policy is feasible and within \(\varepsilon\) of the finite optimum value. Report an empty feasible set separately.

\textbf{Welfare and accounting.} Normative utility, interpersonal normalization, social weights, numeraire, discounting and the use of public receipts are primitives. Behavioral utility need not coincide with the welfare criterion. Household welfare includes ownership income and rents once; adding profits again to compensating variation generally double-counts them. A monetary equivalent is defined by an explicit transfer and price convention, with monotonicity and range conditions. Average effects, mechanism contributions, transfers and welfare losses are different targets. Infinite sums require convergence and dynamic feasible plans require terminal or no-Ponzi conditions.

\textbf{Source versions.} A working-paper date, revision date and journal publication year are distinguished. A DOI for a journal version is not automatically the DOI of an earlier manuscript. Locators refer to the stated version and printed pages where specified. A metadata-only check does not verify a claim about a full-text conclusion. Every entry's source subsection narrows the provenance claim accordingly.
% END ECONOMICS STATEMENT
\end{document}
